\documentclass[11pt]{article}
\usepackage[a4paper,margin=1in]{geometry}
\usepackage{fontspec}
\usepackage[boldfont=false]{xeCJK}
\setCJKmainfont{FandolSong-Regular.otf}
\usepackage{amsmath}
\usepackage{amssymb}
\usepackage{array}
\usepackage{booktabs}
\usepackage{graphicx}
\usepackage{adjustbox}
\usepackage{enumitem}
\setlist[itemize]{topsep=2pt,partopsep=0pt,parsep=0pt,itemsep=2pt,leftmargin=1.4em}
\setlist[enumerate]{topsep=2pt,partopsep=0pt,parsep=0pt,itemsep=2pt,leftmargin=1.6em}
\usepackage{parskip}
\usepackage{fvextra}
\fvset{breaklines=true,breakanywhere=true,fontsize=\small}
\setlength{\parindent}{0pt}

\begin{document}

\begin{center}
  {\LARGE\bfseries Motion, forces and energy}\\[0.35em]
  {\large Igcse Physics}
\end{center}
\vspace{0.6em}

\section*{Measurement 测量}

\subsection*{Length and volume}

Use a \textbf{ruler} 直尺 to measure length. Read the scale with your eye straight in front of the mark to avoid a reading error.

Use a \textbf{measuring cylinder} 量筒 to measure the volume of a liquid. Read the scale at the bottom of the curved surface (the \textbf{meniscus} 弯月面), with your eye level with it.

\subsection*{Measuring small amounts}

A single small length or a short time is hard to measure well. The trick is to \textbf{measure many and divide} 测多个再相除:

\begin{itemize}
\item To find the thickness of one page, measure 100 pages and divide by 100.

\item To find the time for one swing of a \textbf{pendulum} 摆, measure the time for 20 swings and divide by 20. One full swing is called the \textbf{period} 周期.

\end{itemize}

This makes the \textbf{uncertainty} 不确定度 (the size of the error) much smaller.

\subsection*{Scalars and vectors}

A \textbf{scalar} 标量 has size (\textbf{magnitude} 大小) only. A \textbf{vector} 矢量 has size \textbf{and} direction.

\begin{itemize}
\item Scalars: distance, speed, time, mass 质量, energy 能量, temperature 温度.

\item Vectors: force 力, weight 重力, velocity 速度, acceleration 加速度, momentum 动量.

\end{itemize}

To add two vectors at right angles (90°), draw them as two sides of a rectangle. The \textbf{resultant} 合矢量 is the diagonal. Find its size with Pythagoras and its direction with trigonometry.

\[
R = \sqrt{a^2 + b^2}
\]

\section*{Motion 运动}

\subsection*{Speed and velocity}

\textbf{Speed} 速率 is the distance travelled per unit time.

\[
v = \frac{s}{t}
\]

\textbf{Velocity} is speed in a stated direction. So velocity is a vector but speed is a scalar.

\[
\text{average speed} = \frac{\text{total distance}}{\text{total time}}
\]

\subsection*{Acceleration}

\textbf{Acceleration} is the change in velocity per unit time.

\[
a = \frac{\Delta v}{\Delta t}
\]

Here $\Delta v$ means "the change in velocity". A \textbf{deceleration} 减速 (slowing down) is a negative acceleration.

\subsection*{Motion graphs}

A \textbf{distance–time graph} 距离-时间图 shows how far an object has gone:

\begin{itemize}
\item A flat (horizontal) line means the object is \textbf{at rest} 静止.

\item A straight slope means constant speed. The \textbf{gradient} 斜率 (steepness) is the speed.

\item A curve that gets steeper means the object is speeding up.

\end{itemize}

A \textbf{speed–time graph} 速度-时间图 shows how fast an object is going:

\begin{itemize}
\item A flat line means constant speed.

\item A straight slope means constant acceleration. The gradient is the acceleration.

\item The \textbf{area under the line} 线下面积 is the distance travelled.

\end{itemize}

\subsection*{Falling objects}

Near the Earth, all objects speed up as they fall at the same rate. This is the \textbf{acceleration of free fall} 自由落体加速度, $g \approx 9.8\ \text{m/s}^2$.

When an object falls through air, \textbf{air resistance} 空气阻力 (a drag force) acts upward. As it speeds up, air resistance grows. When air resistance equals the weight, the resultant force is zero and the object stops speeding up. It then falls at a steady \textbf{terminal velocity} 终极速度.

\section*{Mass and weight}

\textbf{Mass} is the amount of \textbf{matter} 物质 in an object. It is measured in kilograms (kg) and does not change when you move the object.

\textbf{Weight} is the force of gravity on a mass. It is measured in newtons (N). Weight can change: it is smaller on the Moon because the Moon's gravity is weaker.

\textbf{Gravitational field strength} 重力场强度 is the force per unit mass:

\[
g = \frac{W}{m}
\]

This $g$ has the same value as the acceleration of free fall ($\approx 9.8\ \text{N/kg}$). You can compare masses with a \textbf{balance} 天平.

\section*{Density 密度}

\textbf{Density} is the mass per unit volume.

\[
\rho = \frac{m}{V}
\]

The symbol $\rho$ is the Greek letter "rho". The unit is $\text{kg/m}^3$ or $\text{g/cm}^3$.

To find density: measure the mass with a balance, find the volume, then divide.

\begin{itemize}
\item \textbf{Regular solid} 规则固体 (like a box): measure the sides and calculate the volume.

\item \textbf{Irregular solid} 不规则固体 (a strange shape): lower it into water in a measuring cylinder. The rise in water level is its volume. This is the \textbf{displacement method} 排水法.

\end{itemize}

An object \textbf{floats} 漂浮 if its density is less than the density of the liquid. It sinks if its density is greater.

\section*{Forces}

A force is a push or a pull. A force can change the \textbf{shape} 形状, the speed, or the direction of an object.

\subsection*{Stretching (Hooke's law)}

When you hang a load on a spring, it stretches. The stretch is called the \textbf{extension} 伸长量.

On a \textbf{load–extension graph} 载荷-伸长图 the line is straight at first: extension is proportional to load. The point where the line stops being straight is the \textbf{limit of proportionality} 比例极限.

The \textbf{spring constant} 弹簧常数 is the force per unit extension:

\[
k = \frac{F}{x}
\]

A large $k$ means a stiff spring.

\subsection*{Resultant force and Newton's laws}

Add forces on a straight line to get the \textbf{resultant force} 合力 (forces one way are positive, the other way negative).

\begin{itemize}
\item If the resultant force is zero, the object stays at rest \textbf{or} keeps moving in a straight line at constant speed. (Newton's first law.)

\item If the resultant force is not zero, the object accelerates in the direction of the force:

\end{itemize}

\[
F = ma
\]

\subsection*{Friction}

\textbf{Friction} 摩擦力 is the force between two surfaces that touch. It tries to stop motion and makes things heat up. Drag (friction in a liquid or gas, like air resistance) also slows objects down.

\subsection*{Moments — the turning effect}

The \textbf{moment} 力矩 of a force is its turning effect about a \textbf{pivot} 支点.

\[
\text{moment} = \text{force} \times \text{perpendicular distance from the pivot}
\]

The unit is the newton metre (N m).

\textbf{Principle of moments} 力矩原理: when an object is balanced (in \textbf{equilibrium} 平衡),

\[
\text{total clockwise moment} = \text{total anticlockwise moment}
\]

An object is in equilibrium when there is no resultant force \textbf{and} no resultant moment.

\subsection*{Centre of gravity}

The \textbf{centre of gravity} 重心 is the single point where all the weight of an object seems to act.

For a flat shape (\textbf{lamina} 薄片), hang it from a pin and let it settle; draw a vertical line down from the pin using a plumb line. Repeat from another point. The centre of gravity is where the lines cross.

An object is more \textbf{stable} 稳定 when its centre of gravity is low and its base is wide. It tips over when the centre of gravity passes outside the base.

\section*{Momentum}

\textbf{Momentum} is mass times velocity. It is a vector.

\[
p = mv
\]

\textbf{Impulse} 冲量 is the force times the time it acts, and it equals the change in momentum:

\[
\text{impulse} = F\,\Delta t = \Delta(mv)
\]

So the resultant force is the change in momentum per unit time:

\[
F = \frac{\Delta p}{\Delta t}
\]

\textbf{Conservation of momentum} 动量守恒: when objects collide and no outside force acts, the total momentum before equals the total momentum after.

\[
m_1 u_1 + m_2 u_2 = m_1 v_1 + m_2 v_2
\]

(Here $u$ is a velocity before and $v$ is a velocity after.)

\section*{Energy, work and power}

\subsection*{Energy stores}

Energy can be \textbf{stored} 储存 in different ways: kinetic 动能, gravitational potential 重力势能, chemical 化学能, elastic (strain) 弹性势能, nuclear 核能, electrostatic 静电能, and internal (thermal) 内能.

Energy is \textbf{transferred} 转移 between stores by forces (mechanical work), by electric currents, by heating, and by waves (such as light and sound).

\subsection*{Kinetic and potential energy}

\textbf{Kinetic energy} is the energy of a moving object:

\[
E_k = \frac{1}{2}mv^2
\]

The change in \textbf{gravitational potential energy} when an object goes up or down by a height $\Delta h$:

\[
\Delta E_p = mg\,\Delta h
\]

\subsection*{Conservation of energy}

The \textbf{principle of conservation of energy} 能量守恒定律 says energy is never made or destroyed; it only moves between stores. A falling object turns gravitational potential energy into kinetic energy. A \textbf{Sankey diagram} 桑基图 shows how the input energy splits into useful energy and \textbf{wasted} 浪费 energy.

\subsection*{Work}

\textbf{Work done} 做功 equals the energy transferred. When a force moves an object:

\[
W = Fd = \Delta E
\]

The unit of work and energy is the \textbf{joule} 焦耳 (J).

\subsection*{Power}

\textbf{Power} 功率 is the work done (or energy transferred) per unit time.

\[
P = \frac{W}{t} = \frac{\Delta E}{t}
\]

The unit is the \textbf{watt} 瓦特 (W). $1\ \text{W} = 1\ \text{J/s}$.

\subsection*{Efficiency}

\textbf{Efficiency} 效率 tells you how much of the input energy becomes useful energy.

\[
\text{efficiency} = \frac{\text{useful energy output}}{\text{total energy input}} \times 100\%
\]

Efficiency is always less than 100\% because some energy is always wasted (usually as heat).

\section*{Pressure 压强}

\textbf{Pressure} is the force per unit area.

\[
p = \frac{F}{A}
\]

The unit is the \textbf{pascal} 帕斯卡 (Pa). $1\ \text{Pa} = 1\ \text{N/m}^2$.

A small area gives a large pressure (a sharp knife cuts well). A large area gives a small pressure (snowshoes stop you sinking).

\subsection*{Pressure in a liquid}

In a liquid, pressure increases with \textbf{depth} 深度 and with the liquid's density:

\[
\Delta p = \rho g \,\Delta h
\]

This is why a dam is built thicker at the bottom, where the water pressure is greatest. Pressure in a liquid acts in all directions.


\end{document}
